% ============================================================
%  A5 - Präsentation mit Beamer (Startdatei)
%  Kompilieren: pdflatex folien.tex
%  Die Datei verwendet zunächst das Standard-Theme von Beamer.
%  Aufgaben laut README:
%    1. Folien um Overlays (\pause, <1->, \only) erweitern
%    2. Eigenes Theme beamerthemeFH.sty anlegen
%    3. Unten \usetheme{FH} einkommentieren
% ============================================================
\documentclass[aspectratio=169,11pt]{beamer}
\usepackage[T1]{fontenc}
\usepackage[ngerman]{babel}

% \usetheme{FH}              % lädt beamerthemeFH.sty (Schritt 3)

\title{\LaTeX{} Advanced}
\subtitle{Präsentationen mit Beamer}
\author{Max Mustermann}
\date{\today}

\begin{document}

\begin{frame}[plain]
  \titlepage
\end{frame}

\begin{frame}{Schrittweiser Aufbau}
  % TODO: zwischen den Punkten \pause einfügen
  \begin{itemize}
    \item Eine Folie ist eine \texttt{frame}-Umgebung.
    \item \texttt{\textbackslash pause} teilt die Folie in Schritte.
    \item Aus einer Folie werden so mehrere PDF-Seiten.
  \end{itemize}
\end{frame}

\begin{frame}{Overlays gezielt steuern}
  % TODO: \item<1->, \item<2->, \item<3> verwenden und \only<...>{...} ergänzen
  \begin{itemize}
    \item Dieser Punkt ist ab Schritt 1 sichtbar.
    \item Dieser Punkt ist ab Schritt 2 sichtbar.
    \item Dieser Punkt ist nur in Schritt 3 sichtbar.
  \end{itemize}
\end{frame}

\begin{frame}{Blöcke und Spalten}
  % TODO: zwei Spalten mit columns/column, links ein block, rechts eine Formel
  Blöcke heben Inhalte hervor.
\end{frame}

\end{document}
